\XeTeXlinebreaklocale "zh"
\XeTeXlinebreakskip = 0pt plus 1pt
\emergencystretch=3em
\usepackage{fvextra}
\usepackage{newunicodechar}
\newfontfamily\cataloglatinfont{DejaVu Serif}
\newunicodechar{ć}{{\cataloglatinfont\char"0107}}
\newunicodechar{ő}{{\cataloglatinfont\char"0151}}
\newunicodechar{ń}{{\cataloglatinfont\char"0144}}
\newunicodechar{Č}{{\cataloglatinfont\char"010C}}
\newunicodechar{ĭ}{{\cataloglatinfont\char"012D}}
\DefineVerbatimEnvironment{Highlighting}{Verbatim}{breaklines,commandchars=\\\{\}}
\AtBeginDocument{\renewcommand{\contentsname}{目录}}
\AddToHook{cmd/section/before}{\clearpage}
\renewcommand{\maketitle}{%
\begin{titlepage}
\centering
\vspace*{45mm}
{\sffamily\fontsize{38}{50}\selectfont 综合数学\par}
\vspace{20mm}
{\sffamily\fontsize{18}{28}\selectfont 学习\qquad 实践\qquad 成长\par}
\vspace{34mm}
{\fontsize{17}{26}\selectfont Leo Lee\par}
\vspace{8mm}
{\fontsize{12}{20}\selectfont 一个喜欢数学的普通人写的\par}
\vfill
{\small 第一批工作稿\quad ·\quad 全书持续编写与审校中\par}
\vspace{8mm}
{\large 二〇二六年\par}
\end{titlepage}
\begin{titlepage}
\centering
\vspace*{76mm}
{\fontsize{15}{28}\selectfont
献给每一位愿意从问题出发、追问理由的读者。\par
\vspace{10mm}
愿每一个定义都有来处，\par
每一步推理都有依据。\par}
\vfill
\end{titlepage}
}
